\section{Proofs of the main results}\label{sec:proofs}

Figure~\ref{fig:flow} shows the structure of the proof. Lemmas~\ref{lem:premium}--\ref{lem:LI} are
specific to loss ratios; Lemmas~\ref{lem:cashlip}--\ref{lem:rep} reprove the representation of comonotonic
additive functionals in the form needed here.

\begin{figure}[ht]
\centering\small
\fbox{\begin{minipage}{0.92\textwidth}\centering
\textbf{Lemma~\ref{lem:premium}} (implied premium): $\Pi$ from the price point; cash additive; $\Pi(c)=c$\ \ (\axP, \axN)\\[2pt]
$\Downarrow$\\[2pt]
\textbf{Lemma~\ref{lem:implied}} (implied risk measure): $\rho=\Pi\cdot\LR$ extends; Möbius cash lines\ \ (\axCM{} with sure losses)\\[2pt]
$\Downarrow$\\[2pt]
\textbf{Lemma~\ref{lem:comon}}: $\rho,\Pi$ comonotonic additive, also on the break-even set\ \ (\axCM)\\[2pt]
$\Downarrow$\\[2pt]
\textbf{Lemma~\ref{lem:LI}}: $\rho,\Pi$ law invariant\ \ (\axLI)\\[2pt]
$\Downarrow$\\[2pt]
\textbf{Lemma~\ref{lem:cashlip}}: homogeneity, Lipschitz\ \ (\axRC{} or \axRM)\\[2pt]
$\Downarrow$\\[2pt]
\textbf{Lemma~\ref{lem:layer}}: layer formula; $h(p)=\rho(\1_A)$, $g(p)=\Pi(\1_A)$\ \ $\Rightarrow$\ \
\textbf{Lemma~\ref{lem:bv}}: bounded variation\ \ $\Rightarrow$\ \ \textbf{Lemma~\ref{lem:rep}}: $\rho=I_h$, $\Pi=I_g$
\end{minipage}}
\caption{Logical structure of the proof of Theorems~\ref{thm:A} and~\ref{thm:B}.}\label{fig:flow}
\end{figure}

\begin{lemma}[Implied premium]\label{lem:premium}
Assume \axP. Then $\Pi(X+c)=\Pi(X)+c$ for all $X$ and $c$, and $D=\{\Pi\ne0\}$. If also \axN, then
$\Pi(c)=c$ for every constant $c$.
\end{lemma}

\begin{proof}
The cash line of $X+c$ is $d\mapsto\LR(X+c+d)$, the cash line of $X$ shifted by $c$; its price point is
therefore the price point of $X$ minus $c$, which is \eqref{eq:Picash}. $X\in D$ means that $0$ is not the
price point, i.e.\ $\Pi(X)\ne0$. Under \axN{} the cash line of $0$ takes the value $1$ at every $d\neq0$, so
its price point is $0$, $\Pi(0)=0$, and $\Pi(c)=c$. \qedhere
\end{proof}

\begin{lemma}[Implied risk measure]\label{lem:implied}
Assume \axP, \axN, \axCM. If $X\in D$ and $X+c\in D$, then $\rho(X+c)=\rho(X)+c$. Consequently
$\rho(X):=\rho(X+c)-c$, for any $c\ne-\Pi(X)$, defines a cash-additive functional on $L^\infty$ that agrees
with $\Pi\cdot\LR$ on $D$, satisfies $\rho(c)=c$, and
\begin{equation}\label{eq:mobius}
 \LR(X+c)=\frac{\rho(X)+c}{\Pi(X)+c}\qquad\text{for all }c\ne-\Pi(X).
\end{equation}
\end{lemma}

\begin{proof}
Let $X,X+c\in D$ with $c\ne0$. The constant $c$ lies in $D$ (by \axN), is comonotone with $X$, and
$\Pi(X)+\Pi(c)=\Pi(X+c)\ne0$ by Lemma~\ref{lem:premium}. So \axCM{} applies to the pair $(X,c)$ and gives
$\LR(X+c)=(\Pi(X)\LR(X)+c)/(\Pi(X)+c)$, i.e.\ $\rho(X+c)=\Pi(X+c)\LR(X+c)=\rho(X)+c$.

For arbitrary $X$, $X+c\in D$ exactly when $c\ne-\Pi(X)$. If $c,c'$ are two such values, then
$X+c'=(X+c)+(c'-c)$ with both $X+c$ and $X+c'$ in $D$, so $\rho(X+c')=\rho(X+c)+(c'-c)$; hence
$\rho(X+c)-c$ does not depend on $c$, and for $X\in D$ (take $c=0$) it equals $\Pi(X)\LR(X)$. Cash additivity
of the extension is immediate. For $c\neq-\Pi(X)$,
$\LR(X+c)=\rho(X+c)/\Pi(X+c)=(\rho(X)+c)/(\Pi(X)+c)$. Finally $\rho(c)=\Pi(c)\LR(c)=c$ for $c\ne0$, and
$\rho(0)=\rho(1)-1=0$. \qedhere
\end{proof}

\begin{proof}[Proof of Proposition~\ref{prop:readings}]
(a) is \eqref{eq:mobius}, and (b) restates it. For (c), $(c+\Pi(X))\LR(X+c)=\rho(X)+c\to\rho(X)-\Pi(X)$ as
$c\to-\Pi(X)$. For (d), the right-hand side of \eqref{eq:mobius} vanishes only at $c=-\rho(X)$, which differs
from the price point when $\rho(X)\ne\Pi(X)$; when $\rho(X)=\Pi(X)$ it equals $1$ for every $c\ne-\Pi(X)$.
Uniqueness of the extension is clear, since any cash-additive extension must satisfy (b). \qedhere
\end{proof}

\begin{lemma}[Comonotonic additivity]\label{lem:comon}
Assume \axP{} and \axN.
\begin{enumerate}[label=(\alph*),nosep]
\item If \axCM{} holds, then $\rho$ and $\Pi$ are comonotonic additive.
\item Conversely, if $\Pi$ is comonotonic additive and there is a comonotonic additive $\tilde\rho$ with
  $\tilde\rho=\Pi\cdot\LR$ on $D$, then \axCM{} holds and the implied risk measure equals $\tilde\rho$.
\end{enumerate}
\end{lemma}

\begin{proof}
(a) Let $X,Y$ be comonotone, $Z=X+Y$, $r=\rho(X)+\rho(Y)$ and $q=\Pi(X)+\Pi(Y)$. Fix $s\neq-q$ and choose
$u\notin\{-\Pi(X),\,s+\Pi(Y)\}$.
\begin{enumerate}[label=(\arabic*),nosep]
\item $X+u$ and $Y+s-u$ are comonotone, and by Lemma~\ref{lem:premium} their premiums $\Pi(X)+u$ and
  $\Pi(Y)+s-u$ are non-zero. So both lie in $D$, and their premiums sum to $q+s\neq0$.
\item Apply \axCM{} to this pair and use $\rho=\Pi\cdot\LR$ on $D$ with Lemma~\ref{lem:implied}:
  \[
  \LR(Z+s)=\frac{(\rho(X)+u)+(\rho(Y)+s-u)}{q+s}=\frac{r+s}{q+s}\in\R .
  \]
\item Thus the cash line of $Z$ is finite at every $s\neq-q$. By \axP{} it has exactly one non-finite
  point, so that point is $-q$, and $\Pi(Z)=q$.
\item For $s\ne-q$, Lemma~\ref{lem:implied} gives $\rho(Z)+s=\rho(Z+s)=\Pi(Z+s)\LR(Z+s)=r+s$, so $\rho(Z)=r$.
\end{enumerate}
The argument never asks whether $X$, $Y$ or $Z$ break even: \axP{} locates the price point even when the
cash line equals $1$ everywhere else.

(b) Let $X,Y\in D$ be comonotone with $q=\Pi(X)+\Pi(Y)\ne0$. Then $\Pi(X+Y)=q\ne0$, so $X+Y\in D$, and
$\LR(X+Y)=\tilde\rho(X+Y)/q=(\tilde\rho(X)+\tilde\rho(Y))/q=(\Pi(X)\LR(X)+\Pi(Y)\LR(Y))/q$. Both $\tilde\rho$ and
the implied risk measure are cash additive (the former by additivity with constants and $\tilde\rho(c)=\Pi(c)\LR(c)=c$)
and agree on $D$, hence everywhere. \qedhere
\end{proof}

\begin{lemma}[Law invariance of the implied pair]\label{lem:LI}
Assume \axP, \axN, \axCM{} and \axLI. Then $\rho$ and $\Pi$ are law invariant.
\end{lemma}

\begin{proof}
If $X\eqd Y$ then $X+c\eqd Y+c$ for every $c$, so by \axLI{} the cash lines of $X$ and $Y$ have the same finite
values and the same non-finite locations. Hence they have the same price point, $\Pi(X)=\Pi(Y)$, and by Proposition~\ref{prop:readings}(b) $\rho(X)=\rho(Y)$. \qedhere
\end{proof}

\begin{lemma}[Homogeneity and Lipschitz continuity]\label{lem:cashlip}
Assume \axP, \axN, \axCM, and let $F\in\{\rho,\Pi\}$. If \axRC{} or \axRM{} holds, $F$ is positively homogeneous
and Lipschitz; under \axRM{} it is $1$-Lipschitz.
\end{lemma}

\begin{proof}
First, the axioms transfer from $\rho_D$ to $\rho$. Given $X,Y$, choose $c$ with $|c|$ arbitrarily small and
$X+c,Y+c\in D$ (only two values are excluded). Then $\rho(X)-\rho(Y)=\rho_D(X+c)-\rho_D(Y+c)$ with
$\|(X+c)-(Y+c)\|=\|X-Y\|$, and $X\le Y$ implies $X+c\le Y+c$; so $\rho$ is uniformly continuous with the modulus
of $\rho_D$ under \axRC, and monotone under \axRM.
Since $X$ is comonotone with itself, $F(nX)=nF(X)$ by Lemma~\ref{lem:comon}, so $F(qX)=qF(X)$ for rational $q>0$.
\emph{Under \axRM:} $X\le Y+\|X-Y\|$ gives $F(X)\le F(Y)+\|X-Y\|$ by monotonicity and cash additivity, so $F$ is
$1$-Lipschitz, and homogeneity extends to real $\lambda>0$. \emph{Under \axRC:} continuity gives
$F(\lambda X)=\lambda F(X)$ for real $\lambda>0$. Choose $\delta>0$ with $|F(X')-F(Y')|\le1$ whenever
$\|X'-Y'\|\le\delta$, and apply it to $tX,tY$ with $t=\delta/\|X-Y\|$; this gives
$|F(X)-F(Y)|\le\|X-Y\|/\delta$. \qedhere
\end{proof}

\begin{lemma}[Layer formula]\label{lem:layer}
Assume \axP, \axN, \axCM, \axLI{} and \axRC{} (or \axRM). Define, for $P(A)=p$, $h(p):=\rho(\1_A)$ and
$g(p):=\Pi(\1_A)$. These are well defined by Lemma~\ref{lem:LI}, defined for every $p\in[0,1]$ because the
space is atomless, and $h(0)=g(0)=0$, $h(1)=g(1)=1$. For every chain $C_1\supset\dots\supset C_m$, $x_0\in\R$
and $a_i\ge0$,
\begin{equation}\label{eq:layer}
 \rho\Big(x_0+\sum_i a_i\1_{C_i}\Big)=x_0+\sum_i a_i\,h(P(C_i)),
\end{equation}
and likewise for $\Pi$ with $g$.
\end{lemma}

\begin{proof}
The partial sum $S_{j-1}=\sum_{i<j}a_i\1_{C_i}$ attains its maximum on $C_j$, so it is comonotone with
$\1_{C_j}$. Apply Lemma~\ref{lem:comon} and cash additivity successively, with
$\rho(a\1_C)=a\rho(\1_C)=a\,h(P(C))$ by Lemma~\ref{lem:cashlip}. The values at $0$ and $1$ come from
$\rho(0)=\Pi(0)=0$ and $\rho(1)=\Pi(1)=1$. \qedhere
\end{proof}

\begin{lemma}[Bounded variation]\label{lem:bv}
Under the assumptions of Lemma~\ref{lem:layer}, $h,g\in\Hset_1$. Under \axRM{} they are non-decreasing.
\end{lemma}

\begin{proof}
\emph{Under \axRM.} For $p\le q$ choose nested $A\subset B$ with $P(A)=p$, $P(B)=q$; then $\1_A\le\1_B$ gives
$h(p)\le h(q)$, and likewise for $g$.

\emph{Under \axRC.} Let $L$ be a Lipschitz constant of $\rho$. Take $0=p_0<\dots<p_n=1$ and a uniform random
variable $U$. Put $I_j=\{1-p_j<U\le1-p_{j-1}\}$ and $\delta_j=h(p_j)-h(p_{j-1})$. For $\sigma\in\{0,1\}^n$ let
$W=\sum_j2(n-j)\1_{I_j}$ and $Y_\sigma=W+\sum_j\sigma_j\1_{I_j}$. The values $2(n-j)+\sigma_j$ are strictly
decreasing in $j$, so $W$ and $Y_\sigma$ have the same upper sets $\{U>1-p_j\}$, of probability $p_j$. By
\eqref{eq:layer} and summation by parts, $\rho\big(\sum_jv_j\1_{I_j}\big)=\sum_jv_j\delta_j$ for every
non-increasing $v\ge0$. Hence $|\sum_j\sigma_j\delta_j|=|\rho(Y_\sigma)-\rho(W)|\le L\|Y_\sigma-W\|\le L$.
Choosing $\sigma_j=\1\{\delta_j>0\}$ and then $\sigma_j=\1\{\delta_j<0\}$ gives $\sum_j|\delta_j|\le2L$. \qedhere
\end{proof}

\begin{lemma}[Representation]\label{lem:rep}
Under the assumptions of Lemma~\ref{lem:layer}, $\rho=I_h$ and $\Pi=I_g$.
\end{lemma}

\begin{proof}
For a simple $X=x_0+\sum_ia_i\1_{C_i}$ with $a_i\ge0$ and $C_i=\{X\ge x_i\}$, a direct computation gives
$I_h(X)=x_0+\sum_ia_ih(P(C_i))=\rho(X)$. Simple positions are uniformly dense in $L^\infty$, and $\rho$ and
$I_h$ are both Lipschitz (Lemma~\ref{lem:cashlip} and $|I_h(X)-I_h(Y)|\le\TV(h)\|X-Y\|$), so $\rho=I_h$.
Likewise $\Pi=I_g$. \qedhere
\end{proof}

\begin{proof}[Proof of Theorems~\ref{thm:A} and~\ref{thm:B}]
\emph{(i)$\Rightarrow$(ii).} Lemmas~\ref{lem:premium}--\ref{lem:rep} give $\Pi=I_g$ and $\rho=I_h$ with
$h,g\in\Hset_1$, non-decreasing under \axRM. On $D=\{\Pi\ne0\}=\{I_g\ne0\}$, $\LR=\rho/\Pi=I_h/I_g$, and off $D$
the loss ratio is not finite. Under \ax{S}, at a book with $I_g(X)=0$ the label is $\star$ exactly when the
residue $\rho(X)=I_h(X)$ vanishes, which is the value of $\bracket{I_h(X)}{0}$.

\emph{(ii)$\Rightarrow$(i).} Let $\LR$ be as in (ii). Since $I_g(X+c)=I_g(X)+c$, the cash line of $X$ is non-finite
exactly at $c=-I_g(X)$, so \axP{} holds and $\Pi=I_g$. \axN{} holds because $I_h(c)=I_g(c)=c$. Next, $I_k$ is
comonotonic additive for $k\in\Hset$: for comonotone $X,Y$, any two upper sets $\{X\ge a\}$, $\{Y\ge b\}$ are almost
surely nested (otherwise comonotonicity fails on a $P\otimes P$-positive set); for simple comonotone $X,Y$ the
union of their upper sets is therefore a chain, giving $X=x_0+\sum a_i\1_{C_i}$, $Y=y_0+\sum b_i\1_{C_i}$ with
$a_i,b_i\ge0$, and additivity follows from the layer computation; for general $X,Y$ approximate by the comonotone
simple pairs $\lfloor nX\rfloor/n$, $\lfloor nY\rfloor/n$ and use the Lipschitz bound. Lemma~\ref{lem:comon}(b) with
$\tilde\rho=I_h$ gives \axCM{} and $\rho=I_h$. \axLI{} holds because $I_k$ depends on $X$ only through
$x\mapsto P(X>x)$, \axRC{} because $I_h,I_g$ are Lipschitz, and \axRM{} when $h,g$ are non-decreasing. \qedhere
\end{proof}

\begin{proof}[Proof of Propositions~\ref{prop:unique} and~\ref{prop:derived}]
$g(p)=\Pi(\1_A)$ is minus the price point of $\1_A$, and $h(p)=\rho(\1_A)$ is given by
Proposition~\ref{prop:readings}(b); both are determined by $\LR$, including when $h=g$. The identities for $D$ and
$N$ follow from $\Pi=I_g$, $\rho=I_h$. Monotonicity of $\rho$ is equivalent to monotonicity of $h$ because $I_h$ is
monotone if and only if $h$ is non-decreasing. Proposition~\ref{prop:derived} collects
Lemmas~\ref{lem:premium}--\ref{lem:comon} and~\ref{lem:cashlip}, the Lipschitz bound for $I_k$, and
$\LR(\lambda X)=\lambda\rho(X)/(\lambda\Pi(X))=\LR(X)$ on $D$ (and $\lambda D=D$). \qedhere
\end{proof}
